\subsection{Extended fragment and front-end}\label{sec:impl-fragment}
We extend \textsc{TensorGuard}'s analyzer with four module families
implemented under \texttt{src/v5/} (no edits to legacy
\texttt{src/}). The \emph{symbolic config} environment binds
\texttt{self.config.<attr>} to fresh \texttt{SymInt}s so that
HuggingFace-style modules whose dimensions come from a config object
can be checked once and reused across instantiations
(\texttt{src/v5/symbolic\_config.py}). The \emph{QKV unpacking}
front-end recognises the four idiomatic Q/K/V splits
(\texttt{split}, \texttt{chunk}, \texttt{view+unbind}, and
\texttt{einops.rearrange} with the canonical
\verb|(three h d) -> three b h t d| pattern) and elaborates each
into three correlated tensor types
(\texttt{src/v5/qkv\_unpacking.py}). The \emph{negative-one
reshape} elaborator emits a Z3 divisibility constraint
$P \bmod Q = 0$ where $P$ is the symbolic input product and $Q$ the
product of the explicit reshape dimensions
(\texttt{src/v5/reshape\_neg1.py}). Finally, transfer rules for
\texttt{F.scaled\_dot\_product\_attention},
\texttt{nn.MultiheadAttention}, \texttt{nn.LayerNorm}, and
\texttt{nn.RMSNorm} register into the analyzer's
\texttt{TORCH\_SHAPE\_OPS} dictionary
(\texttt{src/v5/attention\_norms.py}). \emph{74/74}
unit tests pass on \texttt{python3.11} +
\texttt{torch~2.9.1}.
The legacy benchmark targets do not abstain in the first place,
so the v5 contribution lives in the new APIs (callable directly on
realistic Transformer modules) rather than in moving the legacy
verdict counters; we report this honestly in
\texttt{experiments\_v5/track\_C\_coverage.json}.

